\SourceLeaf{272}
\Article{5}
\indent Ich begleite die Vorschriften dieser ganzen Abhandlung mit einem numerischen Beispiel, das mit größter Genauigkeit durchgerechnet ist. Anderen, die ihre Messungen nach der hier vorgetragenen Methode auswerten wollen, kann es entweder als Rechnungsmuster für die Erstellung der erforderlichen Hilfstafeln oder unmittelbar als Hilfsmittel für einen großen Teil der gemäßigten Zone dienen. In den meisten Fällen wird man sich übrigens mit einer \textit{viel} geringeren Genauigkeit begnügen können.

Als Normalbreite wähle ich \(52^0\ 40'\), was ungefähr dem mittleren Breitenkreis des Königreichs Hannover entspricht. Da aber ein runder Wert für die Normalbreite auf der Kugel in mancher Hinsicht vorteilhafter ist als einer für die Normalbreite auf dem Ellipsoid, setze ich
\[
Q=52^0\ 40'\ 0''
\]
\indent Ich führe die Rechnung nach den neuesten von \textsc{Bessel} aus den Gradmessungen abgeleiteten Erddimensionen aus (Astronomische Nachrichten, Band 19, S. 116). Nimmt man die Toise als Einheit, so gilt
\[
\begin{aligned}
\log a&=6{,}5148235337\\
\log\cos\varphi&=9{,}9985458202
\end{aligned}
\]
Daraus folgt mithilfe der zehnstelligen Logarithmen
\[
\begin{aligned}
\varphi&=4^0\ 41'\ 9''98262\\
\log e&=8{,}9122052079\\
\zeta&=1^0\ 43'\ 26''80402\\
\eta&=2\quad15\quad42\ 34083\\
P&=52\quad42\quad2{,}53251\\
\log\alpha&=0{,}0001966553\\
\theta&=3^0\ 43'\ 34''24669\\
\log\frac{1}{k}&=0{,}0016708804\\
\log A&=6{,}5152074703
\end{aligned}
\]
Nimmt man das französische gesetzliche Meter als Einheit, so ergibt sich
\[
\log A=6{,}8050274003
\]
